\section{Interpretação e resultado}

\begin{frame}[t,label=cd-resultado]{6. Interpretação e proposta de ação}
\begin{tabularx}{\linewidth}{@{}lXrrr@{}}
\toprule
\textbf{Curso} & \textbf{Disponíveis} & \textbf{Mediana} & \textbf{Desvio padrão} & \textbf{Ausentes} \\
\midrule
Computação & 6 & 7,85 & 1,29 & 0,0\% \\
Estatística & 5 & 7,80 & 1,38 & 16,7\% \\
Matemática & 5 & 7,50 & 1,25 & 16,7\% \\
\bottomrule
\end{tabularx}
\smallskip

\textbf{Análise:} medianas e desvios padrão próximos entre os cursos.

\textbf{Ação proposta:} conferir as duas notas pendentes com os responsáveis.

\textbf{Avaliação:} buscar 0\% de notas ausentes e repetir a análise.

\begin{EMASCnote}
A base é pequena e fictícia. As medidas descrevem as notas disponíveis, não a qualidade dos cursos.
\end{EMASCnote}
\end{frame}

\begin{frame}[fragile,t]{6.1. Junção de tabelas}
O método \texttt{merge()} combina tabelas a partir de uma coluna em comum. Vamos juntar as medidas das notas ao percentual de notas ausentes pelo curso.
\EMASCcodelabel{Exemplo prático}
\begin{lstlisting}
ausencias = cobertura[["pct_sem_nota"]].reset_index()
resultado = pd.merge(resumo.reset_index(), ausencias, on="curso")
print(resultado.round(2))
\end{lstlisting}
\begin{EMASCnote}
\texttt{reset\_index()} transforma o índice dos cursos em uma coluna. \texttt{on="curso"} indica a coluna usada para relacionar as tabelas.
\end{EMASCnote}
\end{frame}

\begin{frame}[fragile,t]{6.2. Exportação dos resultados}
A tabela final pode ser exportada para utilização em relatórios ou novas análises.
\EMASCcodelabel{Exemplo prático}
\begin{lstlisting}
resultado.round(2).to_csv(
    "resumo_por_curso.csv",
    index=False
)
print("Arquivo resumo_por_curso.csv gerado.")
\end{lstlisting}
\begin{EMASCnote}
\texttt{to\_csv()} grava a tabela em um arquivo CSV. \texttt{index=False} impede que o índice numérico seja gravado como uma coluna adicional.
\end{EMASCnote}
\end{frame}
